\begingroup
\setlength{\tabcolsep}{4.5pt}
\renewcommand{\arraystretch}{1.28}
\begin{longtable}{@{}>{\raggedright\arraybackslash}p{.14\textwidth}>{\raggedright\arraybackslash}p{.33\textwidth}>{\raggedright\arraybackslash}p{.25\textwidth}>{\raggedright\arraybackslash}p{.22\textwidth}@{}}
\caption{Realización de compuestos, sectores topológicos y estados virtuales.}\label{tab:masas-realizacion-no-escalar}\\
\toprule
\textbf{Clase} & \textbf{Objeto HMT} & \textbf{Mapa de realización} & \textbf{Conclusión y condiciones} \\
\midrule
\endfirsthead
\toprule
\textbf{Clase} & \textbf{Objeto HMT} & \textbf{Mapa de realización} & \textbf{Conclusión y condiciones} \\
\midrule
\endhead
\midrule
\multicolumn{4}{r}{\footnotesize Continúa en la página siguiente.}\\
\endfoot
\bottomrule
\endlastfoot
mesones & Registro doble, frontera de enlace y firma compuesta. & Registro compuesto y ligadura $\to$ carta escalar. & Composición de registros; la ligadura precede a la carta escalar. \\
bariones & Registro triple, acarreo y firma compuesta. & Registro compuesto y ligadura $\to$ carta escalar. & Composición de registros; la ligadura precede a la carta escalar. \\
Fibonacci $(1,\tau)$ & Anillo de fusión Fibonacci: $\tau\otimes\tau=1+\tau$. & Fusión y trenzado $\to$ observable físico del sector. & Codominio exacto de fusión y trenzado; el contraste usa sus invariantes propios. \\
Ising $(1,\psi,\sigma)$; realización Majorana separada & Anillo de fusión de Ising: $(1,\psi,\sigma)$. & Fusión y trenzado $\to$ observable físico del sector. & Codominio exacto de fusión y trenzado; el contraste usa sus invariantes propios. \\
sectores $\mathbb Z/6$ & Carácter de trenzado sobre $\mathbb Z/6$. & Fusión y trenzado $\to$ observable físico del sector. & Codominio exacto de fusión y trenzado; el contraste usa sus invariantes propios. \\
secciones virtuales & Sistema inverso de extensiones y horizonte de persistencia. & Prolongación registrada $\to$ horizonte físico de persistencia. & Objeto exacto de prolongación y obstrucción; su salida es el horizonte de persistencia. \\
\end{longtable}
\endgroup
